\documentclass[12pt,a4paper]{article}
\usepackage[margin=25mm, headheight=15pt, headsep=10pt]{geometry}
\usepackage{fontspec}
\usepackage[table]{xcolor} % Note: table option is required for rowcolors
\usepackage{titlesec}
\usepackage{tocloft}
\usepackage{fancyhdr}
\usepackage{booktabs}
\usepackage{tcolorbox}
\tcbuselibrary{skins, breakable}
\usepackage[colorlinks=true, linkcolor=emerald, urlcolor=emerald, bookmarks=true]{hyperref}

\definecolor{emerald}{RGB}{0,102,51}
\definecolor{darkred}{RGB}{139,0,0}
\definecolor{lightgray}{RGB}{245,245,245}

\usepackage[nil]{babel}
\babelprovide[import, main]{bengali}
\babelprovide[import]{arabic}
\babelfont{rm}[Renderer=HarfBuzz,Scale=1.0]{Noto Serif Bengali}
\babelfont{sf}[Renderer=HarfBuzz]{Noto Sans Bengali}
\babelfont{tt}[Renderer=HarfBuzz]{Noto Sans Bengali}
\babelfont[arabic]{rm}[Renderer=HarfBuzz,Scale=1.1]{Noto Naskh Arabic}

% Section styling
\titleformat{\section}{\Large\bfseries\color{emerald}}{\thesection.}{0.5em}{}
\titleformat{\subsection}{\large\bfseries\color{emerald!85!black}}{\thesubsection.}{0.5em}{}
\titleformat{\subsubsection}{\normalsize\bfseries\color{emerald!70!black}}{\thesubsubsection.}{0.5em}{}
\titlespacing*{\section}{0pt}{18pt}{8pt}

% Fancy Header/Footer setup
\pagestyle{fancy}
\fancyhf{} % clear all header and footer fields
\fancyhead[L]{\color{emerald}\small\textbf{\nouppercase{\leftmark}}}
\fancyfoot[L]{\scriptsize\color{black!50}{\lat{AI}}-সংকলিত, আলেম-যাচাইকৃত নয়}
\fancyfoot[C]{\color{emerald!70!black}\thepage}
\renewcommand{\headrulewidth}{0.4pt}
\renewcommand{\headrule}{\hbox to\headwidth{\color{emerald}\leaders\hrule height \headrulewidth\hfill}}
\renewcommand{\footrulewidth}{0pt}

% Custom boxes
% 1. Ayat/Hadith Box
\newtcolorbox{ayatbox}[1][]{
    enhanced, breakable, 
    colback=emerald!5, 
    colframe=emerald, 
    boxrule=0.5pt, 
    arc=3pt, 
    left=10pt, right=10pt, top=10pt, bottom=10pt,
    #1
}

% 2. Quote/Translation Box
\newtcolorbox{quotebox}[1][]{
    enhanced, breakable,
    frame hidden,
    colback=lightgray,
    borderline west={3pt}{0pt}{emerald},
    left=15pt, right=10pt, top=10pt, bottom=10pt,
    sharp corners,
    #1
}

% 3. AI-generated notice box
\newtcolorbox{warnbox}[1][]{
    enhanced, breakable,
    colback=red!4,
    colframe=darkred,
    boxrule=0.8pt,
    arc=3pt,
    left=10pt, right=10pt, top=10pt, bottom=10pt,
    #1
}

% Commands
\newcommand{\arab}[1]{\foreignlanguage{arabic}{#1}}
\newcommand{\kib}[1]{\textcolor{emerald}{\textbf{#1}}}

% Latin (English) fallback: Noto Serif Bengali has NO Latin glyphs
\newfontfamily\latfont[Renderer=HarfBuzz]{Noto Serif}
\newcommand{\lat}[1]{{\latfont #1}}

\setlength{\parskip}{5pt}
\setlength{\parindent}{0pt}

% Fix for missing bullet in Noto Serif Bengali
\renewcommand{\labelitemi}{$\bullet$}



\begin{document}
\begin{center}{\Large\bfseries\color{emerald} কুরআনের আয়াত — অহংকার (কিবর) ও আত্মদাম্ভিকতা (উজব)}\end{center}
\vspace{1em}
\section*{কুরআনের আয়াত — অহংকার (কিবর) ও আত্মদাম্ভিকতা (উজব)}

\begin{warnbox}
 \textbf{গুরুত্বপূর্ণ নোটিশ (\lat{Important} \lat{Notice}):} এই কন্টেন্টটি \lat{AI} (কৃত্রিম বুদ্ধিমত্তা) দিয়ে প্রস্তুত ও সংকলিত। \lat{AI} কঠোর সূত্র-মান নীতি মেনে শুধু নির্ভরযোগ্য উৎস থেকে তথ্য সংগ্রহ করেছে — কেবল সহীহ/হাসান হাদিস ও সহীহ সনদের আসার রাখা হয়েছে, দুর্বল সনদ বাদ দেওয়া হয়েছে। তবে এটি কোনো আলেম বা মুহাদ্দিস কর্তৃক যাচাইকৃত নয়।
\end{warnbox}


\begin{quote}
\textbf{পড়ার নিয়ম:} প্রতিটি আয়াতের সাথে থাকবে — \textbf{মূল আরবি} $\rightarrow$ \textbf{বাংলা অর্থ (\lat{pure})} $\rightarrow$ \textbf{শব্দের ব্যাখ্যা} $\rightarrow$ \textbf{গভীর তাফসীর} $\rightarrow$ \textbf{শিক্ষা ও প্রয়োগ}। \\
তাফসীর সংগ্রহ করা হয়েছে প্রামাণ্য গ্রন্থ থেকে: তাফসীর ইবনে কাসীর, জামিউল বায়ান (তাবারী), তাফসীরুল কুরতুবী, তাফসীর আস-সা'দী, মাআলিমুত তানযীল (বাগাবী), মাআরিফুল কুরআন (মুফতী মুহাম্মদ শাফী উসমানী), ফী যিলালিল কুরআন (সাইয়িদ কুতুব)।
\end{quote}


\medskip\hrule\medskip

\subsection*{আয়াত ১ — সূরা লুকমান (৩১) — আয়াত ১৮}

\subsubsection*{১. সূত্র কার্ড}

\begin{table}[h]\centering\small
\begin{tabular}{ll}
বিষয় & বিবরণ \\
\hline
\textbf{সূরা} & লুকমান (৩১) — মক্কী \\
\textbf{আয়াত} & ১৮ \\
\textbf{পারা} & ২১ \\
\textbf{মূল বিষয়} & অহংকারে মুখ ফেরানো ও দম্ভভরে হাঁটা নিষিদ্ধ \\
\textbf{বক্তা} & লুকমান (আঃ) — তাঁর ছেলেকে উপদেশ দিচ্ছেন \\
\hline
\end{tabular}
\end{table}

\subsubsection*{২. প্রসঙ্গ ও প্রেক্ষাপট (\lat{Context})}

সূরায় লুকমান (আঃ)-এর জীবনের একটি দীর্ঘ \textbf{\lat{upbringing} (লালন-পালন)} অধ্যায় আছে। আয়াত ১৩ থেকে ১৯ পর্যন্ত তিনি ছেলেকে একে একে আদর্শের পাঠ দিচ্ছেন: তাওহীদ, পিতামাতার হক, নামাজ, \textbf{\lat{accountability} (দায়বদ্ধতা)}। এর ধারাবাহিকতায় আয়াত ১৮-এ তিনি ছেলেকে চরিত্রগত \textbf{\lat{moral} \lat{disease} (নৈতিক রোগ)} — অহংকার ও দম্ভ — থেকে বিরত রাখছেন। এটি শুধু ছেলেকে নয়, সমগ্র মানবজাতির জন্য একটি চিরন্তন \textbf{\lat{guideline} (নির্দেশিকা)}।

\subsubsection*{৩. মূল আরবি}
\begin{center}\arab{وَلَا تُصَعِّرْ خَدَّكَ لِلنَّاسِ وَلَا تَمْشِ فِي الْأَرْضِ مَرَحًا ۖ إِنَّ اللَّهَ لَا يُحِبُّ كُلَّ مُخْتَالٍ فَخُورٍ}\end{center}
\textbf{উচ্চারণ:} ওয়া লা- তুসা'য়ির খাদ্দাকা লিন্না-সি ওয়া লা- তামশি ফিল আরদ্বি মারাহা-, ইন্নাল্লা-হা লা- ইউহিব্বু কুল্লা মুখতা-লিন ফাকুর।

\subsubsection*{৪. বাংলা অর্থ (\lat{pure})}
\textbf{\lat{“}আর তুমি মানুষের প্রতি অহংকারবশে মুখ ফিরিয়ে নিও না, আর পৃথিবীতে দম্ভভরে চলো না। নিশ্চয়ই আল্লাহ কোনো আত্মগর্বী, দাম্ভিক ব্যক্তিকে পছন্দ করেন না।\lat{”}}

\subsubsection*{৫. শব্দের ব্যাখ্যা (\lat{Vocabulary})}

\begin{table}[h]\centering\small
\begin{tabular}{ll}
শব্দ & অর্থ \\
\hline
\textbf{\arab{تُصَعِّرْ}} & মুখ বাঁকানো — ঘৃণা ও অহংকারে কারো দিকে মুখ ফেরানো \\
\textbf{\arab{خَدَّكَ}} & তোমার গাল/মুখ \\
\textbf{\arab{مَرَحًا}} & দম্ভভরে হাঁটা — আত্মম্ভরিতায় ভরা চলন \\
\textbf{\arab{مُخْتَالٍ}} & আত্মগর্বী — নিজেকে বড় ভাবা (উজবের রূপ) \\
\textbf{\arab{فَخُورٍ}} & দাম্ভিক — নিজের প্রশংসা করে বেড়ানো \\
\hline
\end{tabular}
\end{table}

\subsubsection*{৬. গভীর তাফসীর (\lat{Deep} \lat{Tafsir})}

\begin{itemize}
\item \textbf{ইবনে কাসীর:} \lat{“}\arab{تُصَعِّرْ}\lat{”} অর্থাৎ \textbf{\lat{humiliation} (অপমান)} ও অহংকারে মুখ সরানো। লুকমান ছেলেকে দুটি জিনিস থেকে বিরত রাখলেন: (১) মানুষের প্রতি মুখ ফিরানো ও তুচ্ছ করা, (২) অহংকারে হাঁটা। কারণ আল্লাহ দাম্ভিক-গর্বিতদের ভালোবাসেন না। মুখ ফেরানো আর দম্ভে চলা — দুটিই কিবরের বাহ্যিক \textbf{\lat{expression} (প্রকাশ)}। (তাফসীর ইবনে কাসীর, ৩/৪৫৬)
\item \textbf{তাবারী:} আল্লাহ সেই ব্যক্তির নিন্দা করছেন যে অহংকারবশে মানুষকে হেয় জ্ঞান করে, তাদের সাথে কথা বলার সময় \textbf{\lat{arrogance} (অহংকার)} -এর কারণে মুখ ঘুরিয়ে নেয়। (জামিউল বায়ান, ২১/৮৪)
\item \textbf{কুরতুবী:} কিবর যে গুণে শোভা পায় তার একটি উদাহরণ দাঁড়িয়ে হাঁটা ও বুক ফুলিয়ে চলা। কোনো মানুষকে বুক ফুলিয়ে চলতে দেখলে মনে হবে সে পৃথিবীর সবকিছু জিতে নিয়েছে, অথচ সে তো মাটিরই সৃষ্টি। (তাফসীরুল কুরতুবী, ১৪/৭৫)
\item \textbf{সা'দী:} আল্লাহ ঐ ব্যক্তিকে \textbf{\lat{dislike} (অপছন্দ)} করেন যে নিজেকে বড় ভাবে (মুখ ফেরানো, চলনে দম্ভ) এবং নিজের প্রশংসা করে (ফখর)। এটা কিবর ও উজবের \textbf{\lat{combination} (সমষ্টি)} — বাইরের অহংকার (কিবর) আর ভেতরের আত্মতুষ্টি (উজব), দুটোই এখানে একসাথে। (তাফসীর আস-সা'দী, পৃষ্ঠা ৬৫৩)
\item \textbf{মাআরিফুল কুরআন:} লুকমান (আঃ)-এর এই উপদেশ মুসলিম সমাজের চিরন্তন আদব। মুখ ফেরানো ও দম্ভভরা চলা অহংকারের \textbf{\lat{outward} \lat{sign} (বাহ্যিক লক্ষণ)}। নিজের জন্য কিবরকে আল্লাহ নিষিদ্ধ করেছেন, বরং কিবর-কিবরিয়া শুধু আল্লাহরই \textbf{\lat{exclusive} (একান্ত)} গুণ। (মাআরিফুল কুরআন, ৭/২৭৯-২৮০)
\end{itemize}
\subsubsection*{৭. সংশ্লিষ্ট হাদিস ও আয়াত}

\begin{itemize}
\item \textbf{হাদিস (মুসলিম ৯১):} \lat{“}যার হৃদয়ে সরিষার দানার ওজনের কিবর থাকবে, সে জান্নাতে প্রবেশ করবে না।\lat{”}
\item \textbf{আয়াত (ইসরা ১৭:৩৭):} \lat{“}আর তুমি পৃথিবীতে দম্ভভরে চলো না...\lat{”} — একই উপদেশ ভিন্ন সূরায়।
\item \textbf{আয়াত (নিসা ৪:৩৬):} \lat{“}নিশ্চয়ই আল্লাহ দাম্ভিক (মুখতাল) ও আত্মপ্রশংসাকারী (ফাখুর) কাউকে পছন্দ করেন না।\lat{”}
\end{itemize}
\subsubsection*{৮. গভীর শিক্ষা (\lat{Deep} \lat{Lessons})}

\begin{enumerate}
\item \textbf{অহংকারের প্রথম চিহ্ন — মানুষের সাথে আচরণে।} আয়াত শুরু হয়েছে \lat{“}মুখ ফেরানো\lat{”} দিয়ে। অর্থাৎ কিবরের প্রথম শিকার হয় মানুষের \textbf{\lat{relationship} (সম্পর্ক)}। আমরা যখন কারো সাথে কথা বলতে অনীহা প্রকাশ করি, তাকে \textbf{\lat{ignore} (উপেক্ষা)} করি — সেটিই অহংকারের বীজ।
\item \textbf{হাঁটার ভঙ্গিও আমল।} ইসলাম শুধু অন্তর নয়, চলন-বলনেরও \textbf{\lat{discipline} (শৃঙ্খলা)} দেয়। দম্ভে হাঁটা আল্লাহর অপছন্দ।
\item \textbf{কিবর ও উজব একসাথে।} আয়াতে \lat{“}\arab{مُخْتَالٍ}\lat{”} (আত্মগর্বী) ও \lat{“}\arab{فَخُورٍ}\lat{”} (দাম্ভিক) — দুটি শব্দই এসেছে। প্রথমটি ভেতরের আত্মতুষ্টি, দ্বিতীয়টি বাইরের প্রদর্শন। মুমিনের দায়িত্ব দুটিই থেকে বেঁচে থাকা।
\item \textbf{বিনয়ই মানুষের মর্যাদা।} যেখানে অহংকার অপছন্দ, সেখানে বিনয় আল্লাহর কাছে প্রিয় — এটি নবীজির সুন্নাত ও কুরআনের শিক্ষা।
\end{enumerate}
\subsubsection*{৯. জীবনে প্রয়োগ (\lat{Practical} \lat{Application})}

\begin{enumerate}
\item কারো সাথে কথা বলার সময় তার দিকে মুখ রাখুন — ঘৃণায় মুখ ফেরাবেন না।
\item হাঁটার ভঙ্গিতে নজর দিন — ভাব যেন দম্ভ না হয়।
\item নিজের প্রশংসা করে বেড়ানো থেকে বিরত থাকুন — প্রশংসা আল্লাহর কাছে।
\item প্রতিদিন একটি করে অহংকারের চিহ্ন নিজের মধ্যে খুঁজুন ও \textbf{\lat{sincerely} (আন্তরিকভাবে)} সংশোধন করুন।
\end{enumerate}
\medskip\hrule\medskip

\subsection*{আয়াত ২ — সূরা আন-নাহল (১৬) — আয়াত ২৩}

\subsubsection*{১. সূত্র কার্ড}

\begin{table}[h]\centering\small
\begin{tabular}{ll}
বিষয় & বিবরণ \\
\hline
\textbf{সূরা} & আন-নাহল (১৬) — মক্কী \\
\textbf{আয়াত} & ২৩ \\
\textbf{পারা} & ১৪ \\
\textbf{মূল বিষয়} & আল্লাহ অহংকারীদের পছন্দ করেন না \\
\hline
\end{tabular}
\end{table}

\subsubsection*{২. প্রসঙ্গ ও প্রেক্ষাপট (\lat{Context})}

আয়াত ২২-২৩ একই প্রসঙ্গের অংশ। আয়াত ২২-এ বলা হয়েছে — কাফিরদের \textbf{\lat{illogical} (অযৌক্তিক)} প্রার্থনা কেবল মিথ্যার উপর। এরপর আয়াত ২৩-এ আল্লাহ বলছেন — তারা গোপনে ও প্রকাশ্যে যা-ই করুক, সবই আমার জ্ঞানে; আর আমি \lat{“}\arab{مُسْتَكْبِرِين}\lat{”} (অহংকারীদের) পছন্দ করি না। অর্থাৎ অহংকারী — যারা হকের সামনে মাথা নত করে না — তাদের কোনো লুকানো আশ্রয় নেই।

\subsubsection*{৩. মূল আরবি}
\begin{center}\arab{لَا جَرَمَ أَنَّ اللَّهَ يَعْلَمُ مَا يُسِرُّونَ وَمَا يُعْلِنُونَ ۚ إِنَّهُ لَا يُحِبُّ الْمُسْتَكْبِرِينَ}\end{center}
\textbf{উচ্চারণ:} লা- জারা-মা আন্নাল্লা-হা ইয়া'লামু মা- ইউসিররূনা ওয়া মা- ইউ'লিনূনা, ইন্নাহূ লা- ইউহিব্বুল মুসতাকবিরীন।

\subsubsection*{৪. বাংলা অর্থ (\lat{pure})}
\textbf{\lat{“}নিশ্চয়ই আল্লাহ জানেন যা তারা গোপনে করে ও যা প্রকাশ্যে করে। নিশ্চয়ই তিনি অহংকারীদের পছন্দ করেন না।\lat{”}}

\subsubsection*{৫. শব্দের ব্যাখ্যা (\lat{Vocabulary})}

\begin{table}[h]\centering\small
\begin{tabular}{ll}
শব্দ & অর্থ \\
\hline
\textbf{\arab{لَا} \arab{جَرَمَ}} & নিশ্চয়ই — অবশ্যই \\
\textbf{\arab{يُسِرُّونَ}} & গোপনে করে — \lat{secret} \lat{deeds} \\
\textbf{\arab{يُعْلِنُونَ}} & প্রকাশ্যে করে — \lat{open} \lat{deeds} \\
\textbf{\arab{الْمُسْتَكْبِرِينَ}} & অহংকারীরা — যারা হকের সামনে মাথা নত করে না \\
\hline
\end{tabular}
\end{table}

\subsubsection*{৬. গভীর তাফসীর (\lat{Deep} \lat{Tafsir})}

\begin{itemize}
\item \textbf{ইবনে কাসীর:} পূর্বের আয়াতে (১৬:২২) বলা হয়েছে তারা যে মিথ্যা ও শিরক করে, আল্লাহ তা জানেন। এখানে বলা হচ্ছে — তারা গোপনে ও প্রকাশ্যে যা করে সবই আল্লাহর জ্ঞানে আছে; আর আল্লাহ \lat{“}\arab{مُسْتَكْبِرِين}\lat{”}-দের ভালোবাসেন না। এখানে \lat{“}অহংকারী\lat{”} বলতে কে — তা প্রেক্ষাপটে স্পষ্ট: যারা তাওহীদের \textbf{\lat{truth} (সত্য)} জেনেও মেনে নেয় না। (ইবনে কাসীর, ২/৫৭৯)
\item \textbf{বাগাবী:} \lat{“}\arab{يُسِرُّونَ}\lat{”} — কুফরকে গোপন রাখা, \lat{“}\arab{يُعْلِنُونَ}\lat{”} — প্রকাশ্যে কুফর ও অন্যায়। উভয় অবস্থাতেই আল্লাহ তাদের হিসাব নেবেন। অহংকারীর লুকানো চক্রান্তও আল্লাহর জ্ঞান থেকে বাইরে নয়। (মাআলিমুত তানযীল, ৩/৮৫)
\item \textbf{সা'দী:} মুস্তাকবির (অহংকারী) সে, যে সত্যটুকু জেনে তবু মেনে নেয় না, মানুষকে ছোট করে। এরূপ ব্যক্তি আল্লাহর ভালোবাসার \textbf{\lat{outside} (বাইরে)}। অহংকার আল্লাহর ভালোবাসার পথে সবচেয়ে বড় \textbf{\lat{barrier} (বাধা)}। (আস-সা'দী, পৃষ্ঠা ৪৪১)
\end{itemize}
\subsubsection*{৭. সংশ্লিষ্ট হাদিস ও আয়াত}

\begin{itemize}
\item \textbf{হাদিস (মুসলিম ৯১):} \lat{“}কিবর হলো সত্যকে তুচ্ছ করা ও মানুষকে হেয় করা।\lat{”} — মুস্তাকবিরের সংজ্ঞা এখানেই।
\item \textbf{আয়াত (সাফফাত ৩৭:৩৫):} \lat{“}তারা যখন কালেমা শুনত, তখন অহংকার করত।\lat{”} — হক মেনে না নেওয়ার অহংকার।
\item \textbf{আয়াত (আ'রাফ ৭:১৪৬):} অহংকারীদের থেকে আল্লাহ সত্যবোধ কেড়ে নেন।
\end{itemize}
\subsubsection*{৮. গভীর শিক্ষা (\lat{Deep} \lat{Lessons})}

\begin{enumerate}
\item \textbf{অহংকার — জ্ঞানের সাথে ঈমানের সম্পর্ক নষ্ট করে।} মানুষ জ্ঞানী হতে পারে, কিন্তু অহংকারের কারণে সেই জ্ঞান তাকে \textbf{\lat{guidance} (পথপ্রদর্শন)} দেয় না।
\item \textbf{গোপন-প্রকাশ্য সবই হিসাবের মধ্যে।} অহংকারী যতই লুকিয়ে থাকুক, আল্লাহর জ্ঞান থেকে কোনো \textbf{\lat{escape} (পলায়ন)} নেই।
\item \textbf{হকের সামনে বিনয়ই ঈমানের প্রমাণ।} যে সত্য শুনে মাথা নত করে, সে আল্লাহর ভালোবাসার পাত্র।
\item \textbf{ভালোবাসা ও অপছন্দ আল্লাহর কাছে।} মুমিনের লক্ষ্য — আল্লাহর ভালোবাসা অর্জন করা; অহংকার সেই লক্ষ্য থেকে মানুষকে ঠেলে দেয়।
\end{enumerate}
\subsubsection*{৯. জীবনে প্রয়োগ (\lat{Practical} \lat{Application})}

\begin{enumerate}
\item সত্য শুনলে সাথে সাথে মেনে নেওয়ার \textbf{\lat{habit} (অভ্যাস)} গড়ুন — শোনার আগে তর্ক নয়।
\item নিজের লুকানো দোষগুলোতে নজর দিন — অহংকার প্রায়ই গোপন অবস্থায় থাকে।
\item প্রতিদিন একটি সত্য মেনে নিন — যেটা মেনে নিতে মন চায় না; এটাই অহংকার ভাঙার \textbf{\lat{exercise} (ব্যায়াম)}।
\end{enumerate}
\medskip\hrule\medskip

\subsection*{আয়াত ৩ — সূরা গাফির / মুমিন (৪০) — আয়াত ৫৬}

\subsubsection*{১. সূত্র কার্ড}

\begin{table}[h]\centering\small
\begin{tabular}{ll}
বিষয় & বিবরণ \\
\hline
\textbf{সূরা} & গাফির / মুমিন (৪০) — মক্কী \\
\textbf{আয়াত} & ৫৬ \\
\textbf{পারা} & ২৪ \\
\textbf{মূল বিষয়} & বিনা দলিলে তর্কের মূল কারণ কিবর \\
\hline
\end{tabular}
\end{table}

\subsubsection*{২. প্রসঙ্গ ও প্রেক্ষাপট (\lat{Context})}

এই সূরায় কাফিরদের সাথে রাসূল (সা.)-এর \textbf{\lat{dialogue} (কথোপকথন)} ও তাদের যুক্তিহীন বিতর্কের কথা আছে। আয়াত ৫৬-এ আল্লাহ বলছেন — যারা আয়াত নিয়ে বিনা \textbf{\lat{evidence} (প্রমাণ)} তর্ক করে, তাদের বুকে কিছুই নেই, কেবল অহংকার। তারা যে মর্যাদা বা লক্ষ্য চায়, তা কখনো পাবে না। নবীজিকে নির্দেশ দেওয়া হয়েছে — এমন ফাসির মোকাবিলায় আল্লাহর কাছে \textbf{\lat{refuge} (আশ্রয়)} চাইতে।

\subsubsection*{৩. মূল আরবি}
\begin{center}\arab{إِنَّ الَّذِينَ يُجَادِلُونَ فِي آيَاتِ اللَّهِ بِغَيْرِ سُلْطَانٍ أَتَاهُمْ ۙ إِن فِي صُدُورِهِمْ إِلَّا كِبْرٌ مَّا هُم بِبَالِغِيهِ ۚ فَاسْتَعِذْ بِاللَّهِ ۖ إِنَّهُ هُوَ السَّمِيعُ الْبَصِيرُ}\end{center}
\textbf{উচ্চারণ:} ইন্নাল্লাযীনা ইউজাদিলূনা ফী আ-য়া-তিল্লা-হি বিগাইরি সুলত্বা-নিন আতা-হুম, ইন ফী সুদূরুহিম ইল্লা- কিবরুম মা- হুম বিবা-লিগীহি, ফাসতা'ইয বিল্লা-হি, ইন্নাহূ হুওয়াস সামী'উল বাসীর।

\subsubsection*{৪. বাংলা অর্থ (\lat{pure})}
\textbf{\lat{“}নিশ্চয়ই যারা আল্লাহর আয়াত সম্পর্কে তর্ক করে, তাদের কাছে কোনো প্রমাণ না এসে পৌঁছানোর আগেই — তাদের অন্তরে শুধু অহংকারই আছে; তারা তা (তাদের কাঙ্ক্ষিত লক্ষ্য) কখনো অর্জন করতে পারবে না। কাজেই আপনি আল্লাহর কাছে আশ্রয় চান। নিশ্চয়ই তিনি সর্বশ্রোতা, সর্বদর্শী।\lat{”}}

\subsubsection*{৫. শব্দের ব্যাখ্যা (\lat{Vocabulary})}

\begin{table}[h]\centering\small
\begin{tabular}{ll}
শব্দ & অর্থ \\
\hline
\textbf{\arab{يُجَادِلُونَ}} & তর্ক-বিতর্ক করে — \lat{argument} \\
\textbf{\arab{سُلْطَانٍ}} & প্রমাণ — দলিল, \lat{evidence} \\
\textbf{\arab{كِبْرٌ}} & অহংকার — এখান থেকেই কিবর শব্দের উৎপত্তি \\
\textbf{\arab{مَا} \arab{هُم} \arab{بِبَالِغِيهِ}} & তারা তা অর্জন করতে পারবে না \\
\textbf{\arab{فَاسْتَعِذْ}} & আশ্রয় চান \\
\hline
\end{tabular}
\end{table}

\subsubsection*{৬. গভীর তাফসীর (\lat{Deep} \lat{Tafsir})}

\begin{itemize}
\item \textbf{ইবনে কাসীর:} বিনা দলিলে তর্ককারীদের অন্তরে কেবল বড়ত্ব ও অহংকার আছে, তাদের বড়ত্ব কখনো অর্জিত হবে না। বরং এর ফলে আল্লাহ তাদের লাঞ্ছিত করবেন। এখানে \lat{“}\arab{كِبْر}\lat{”} শব্দটি থেকেই অহংকারের (কিবর) নামকরণ; অর্থাৎ \lat{arrogance} হলো এমন একটি গুণ যা মানুষকে হকের সামনে দাঁড় করিয়ে দেয়। (ইবনে কাসীর, ৪/৯২-৯৩)
\item \textbf{তাবারী:} যারা রাসূলের বিরোধিতায় কিবর ও গর্বে অটল থাকে, তারা এই সত্যে পৌঁছাতে পারবে না। রাসূলকে বলা হচ্ছে — এই জালিমদের মোকাবিলায় আল্লাহর শরণ নিন। (তাবারী, ২১/৫৬)
\item \textbf{কুরতুবী:} আয়াতে সেই কিবরের উল্লেখ আছে যার ফলে মানুষ হক মেনে নেয় না — এটাই সবচেয়ে মারাত্মক কিবর। যে অহংকার মানুষকে সত্য গ্রহণ থেকে বিরত রাখে, তা জান্নাতের পথে \textbf{\lat{fatal} (ঘাতক)}। (কুরতুবী, ১৫/৩২৩)
\item \textbf{মাআরিফুল কুরআন:} বিনা প্রমাণে আল্লাহর আয়াত নিয়ে বিতর্কের মূল কারণ কিবর। রাসূলকে নির্দেশ দেওয়া হয়েছে ফাসীর কুফফার থেকে আল্লাহর কাছে আশ্রয় চাইতে। (মাআরিফুল কুরআন, ৮/২২০-২২১)
\end{itemize}
\subsubsection*{৭. সংশ্লিষ্ট হাদিস ও আয়াত}

\begin{itemize}
\item \textbf{হাদিস (মুসলিম ৯১):} \lat{“}যার হৃদয়ে সরিষার দানার কিবর থাকবে, সে জান্নাতে প্রবেশ করবে না।\lat{”}
\item \textbf{আয়াত (বাকারাহ ২:২০৬):} \lat{“}যখন তাকে বলা হয় — আল্লাহকে ভয় করো, তখন গর্ব তাকে পাপের দিকে টেনে নেয়।\lat{”}
\item \textbf{হাদিস (আবু দাউদ ৪০৯২):} সুন্দর লোকের প্রশ্ন — নবীজি বললেন: কিবর সত্যকে তুচ্ছ করা ও মানুষকে হেয় করা।
\end{itemize}
\subsubsection*{৮. গভীর শিক্ষা (\lat{Deep} \lat{Lessons})}

\begin{enumerate}
\item \textbf{যুক্তিহীন তর্ক — অহংকারের মুখোশ।} যে মানুষের কাছে সঠিক \lat{evidence} নেই, অথচ তর্ক করেই যায় — তার পেছনে থাকে অহংকার।
\item \textbf{অহংকারের কাঙ্ক্ষিত লক্ষ্য কখনো অর্জিত হয় না।} যারা অহংকার করে বড় হতে চায়, আল্লাহ তাদের বিপরীত দিকে নামান — এটা আল্লাহর \textbf{\lat{law} (আইন)}।
\item \textbf{আল্লাহর কাছে আশ্রয় — অহংকারীদের মোকাবিলার উপায়।} নবীজিকে বলা হয়েছে \lat{“}\arab{فَاسْتَعِذْ} \arab{بِاللَّهِ}\lat{”} — নিজের শক্তি নয়, আল্লাহর শরণ।
\item \textbf{হকের সামনে বিনয় — মুমিনের পরিচয়।} হক জেনে তার সামনে মাথা নত করাই ঈমান।
\end{enumerate}
\subsubsection*{৯. জীবনে প্রয়োগ (\lat{Practical} \lat{Application})}

\begin{enumerate}
\item কোনো বিষয়ে দলিল না থাকলে তর্ক এড়িয়ে চলুন — \textbf{\lat{debate} (বিতর্ক)} করার আগে পড়াশোনা করুন।
\item যে সত্য মেনে নিতে মন চায় না, সেটা নিয়ে নিজেকে প্রশ্ন করুন: \lat{“}এটা কি অহংকার নয়?\lat{”}
\item কঠিন \textbf{\lat{situation} (পরিস্থিতি)}-এ নিজের শক্তির উপর ভরসা না করে আল্লাহর কাছে আশ্রয় চান।
\end{enumerate}
\medskip\hrule\medskip

\subsection*{আয়াত ৪ — সূরা আয-যুমার (৩৯) — আয়াত ৭২}

\subsubsection*{১. সূত্র কার্ড}

\begin{table}[h]\centering\small
\begin{tabular}{ll}
বিষয় & বিবরণ \\
\hline
\textbf{সূরা} & আয-যুমার (৩৯) — মক্কী \\
\textbf{আয়াত} & ৭২ \\
\textbf{পারা} & ২৪ \\
\textbf{মূল বিষয়} & অহংকারীদের আবাস জাহান্নাম \\
\hline
\end{tabular}
\end{table}

\subsubsection*{২. প্রসঙ্গ ও প্রেক্ষাপট (\lat{Context})}

আয়াত ৭১-এ বলা হয়েছে — কাফিরদের দলে দলে জাহান্নামের দিকে নিয়ে যাওয়া হবে, দরজা খুলে দেওয়া হবে, দারোয়ানরা জিজ্ঞেস করবে — \lat{“}তোমাদের কাছে কি রাসূল আসেনি?\lat{”} তারপর আয়াত ৭২-এ সেই \textbf{\lat{final} \lat{verdict} (চূড়ান্ত রায়)} — তাদের বলা হবে জাহান্নামের দরজা দিয়ে প্রবেশ করো। লক্ষণীয় — এখানে \lat{“}\arab{الْمُتَكَبِّرِينَ}\lat{”} (অহংকারীদের) বলা হয়েছে, কারণ জাহান্নামে যাওয়ার মূল \textbf{\lat{cause} (কারণ)} হলো হকের সামনে অহংকার।

\subsubsection*{৩. মূল আরবি}
\begin{center}\arab{قِيلَ ادْخُلُوا أَبْوَابَ جَهَنَّمَ خَالِدِينَ فِيهَا ۖ فَبِئْسَ مَثْوَى الْمُتَكَبِّرِينَ}\end{center}
\textbf{উচ্চারণ:} ক্বীলাদ খুলূ আবওয়া-বা জাহান্নামা খা-লিদীনা ফীহা-, ফাবি'সা মাছওয়াল মুতাকাব্বিরীন।

\subsubsection*{৪. বাংলা অর্থ (\lat{pure})}
\textbf{\lat{“}(ফেরেশতারা) বলবে: চিরকাল থাকার জন্য জাহান্নামের দরজা দিয়ে প্রবেশ করো। অতএব অহংকারীদের আবাসস্থল কতই না নিকৃষ্ট!\lat{”}}

\subsubsection*{৫. শব্দের ব্যাখ্যা (\lat{Vocabulary})}

\begin{table}[h]\centering\small
\begin{tabular}{ll}
শব্দ & অর্থ \\
\hline
\textbf{\arab{أَبْوَابَ} \arab{جَهَنَّمَ}} & জাহান্নামের দরজাগুলো \\
\textbf{\arab{خَالِدِينَ}} & চিরস্থায়ী — চিরকাল থাকবে \\
\textbf{\arab{مَثْوَى}} & আবাস — \lat{residence} \\
\textbf{\arab{الْمُتَكَبِّرِينَ}} & অহংকারীরা — যারা হকের সামনে অহংকার করেছে \\
\hline
\end{tabular}
\end{table}

\subsubsection*{৬. গভীর তাফসীর (\lat{Deep} \lat{Tafsir})}

\begin{itemize}
\item \textbf{ইবনে কাসীর:} যারা দুনিয়ায় হকের সামনে অহংকার দেখিয়েছে, তাদেরকে বলা হবে জাহান্নামের দরজা দিয়ে প্রবেশ করো। এটাই মুস্তাকবিরিনের প্রতিদান। (ইবনে কাসীর, ৪/৬৪)
\item \textbf{সা'দী:} তাদের অহংকারের জন্যই তাদেরকে জাহান্নামের নিচের স্তরে প্রবেশ করানো হবে। জাহান্নামের দরজা পর্যন্ত তাদেরকে অসম্মানের সাথে ঠেলে নেওয়া হবে। দুনিয়ায় তারা যে \textbf{\lat{respect} (সম্মান)} চেয়েছিল, সেখানে পাবে সম্পূর্ণ বিপরীত। (আস-সা'দী, পৃষ্ঠা ৭১৪)
\item \textbf{মাআরিফুল কুরআন:} এই আয়াতের প্রসঙ্গে সহীহ বুখারিতে আছে — জান্নাত ও জাহান্নাম সৃষ্টির পর উভয়কে জিজ্ঞেস করা হলো; জাহান্নাম বলল, আমার মধ্যে অহংকারীরা আসবে। (মাআরিফুল কুরআন, ৮/১৯৪)
\end{itemize}
\subsubsection*{৭. সংশ্লিষ্ট হাদিস ও আয়াত}

\begin{itemize}
\item \textbf{হাদিস (মুসলিম ২৮৬৫):} \lat{“}জাহান্নাম বলল: আমার মধ্যে অহংকারী ও উদ্ধতরা আসবে।\lat{”}
\item \textbf{আয়াত (আ'রাফ ৭:৪০):} \lat{“}অহংকারীদের জন্য আকাশের দরজা খোলা হবে না।\lat{”}
\item \textbf{হাদিস (মুসলিম ৯১):} \lat{“}যার হৃদয়ে সরিষার দানার কিবর থাকবে, সে জান্নাতে প্রবেশ করবে না।\lat{”}
\end{itemize}
\subsubsection*{৮. গভীর শিক্ষা (\lat{Deep} \lat{Lessons})}

\begin{enumerate}
\item \textbf{অহংকারের চূড়ান্ত পরিণতি — জাহান্নাম।} দুনিয়ায় অহংকারী যতই \textbf{\lat{glory} (গৌরব)} খুঁজুক, আখেরাতে তার ঠিকানা জাহান্নাম।
\item \textbf{প্রবেশের আদেশ — অসম্মানের সাথে।} তাদের বলা হচ্ছে \lat{“}প্রবেশ করো\lat{”} — কিন্তু এটি সম্মান নয়, শাস্তির \textbf{\lat{announcement} (ঘোষণা)}।
\item \textbf{হকের সামনে অহংকার — জাহান্নামের প্রধান কারণ।} আয়াতটি বিশেষভাবে অহংকারীদের উল্লেখ করেছে — কারণ অহংকারই ঈমান গ্রহণে বাধা দেয়।
\item \textbf{চিরস্থায়ীতা।} \lat{“}\arab{خَالِدِينَ}\lat{”} — চিরকাল। অহংকারের শাস্তি সাময়িক নয়; এটা \textbf{\lat{permanent} (স্থায়ী)} বন্দোবস্ত।
\end{enumerate}
\subsubsection*{৯. জীবনে প্রয়োগ (\lat{Practical} \lat{Application})}

\begin{enumerate}
\item নিজেকে প্রশ্ন করুন: \lat{“}আমি কোন সত্যের সামনে অহংকার করি?\lat{”} — সেটিই সংশোধন করুন।
\item আখেরাতের \textbf{\lat{Day} \lat{of} \lat{Judgment} (হিসাবের দিন)} স্মরণে রাখুন — এতে অহংকার কমবে।
\item দুনিয়ার সম্মানের \textbf{\lat{obsession} (মোহ)} ত্যাগ করুন — প্রকৃত সম্মান তাকওয়ায়।
\end{enumerate}
\subsection*{আয়াত ৫ — সূরা বনী ইসরাঈল / আল-ইসরা (১৭) — আয়াত ৩৭}

\subsubsection*{১. সূত্র কার্ড}

\begin{table}[h]\centering\small
\begin{tabular}{ll}
বিষয় & বিবরণ \\
\hline
\textbf{সূরা} & বনী ইসরাঈল / আল-ইসরা (১৭) — মক্কী \\
\textbf{আয়াত} & ৩৭ \\
\textbf{পারা} & ১৫ \\
\textbf{মূল বিষয়} & মানুষের দুর্বলতা — দম্ভ নিরর্থক \\
\hline
\end{tabular}
\end{table}

\subsubsection*{২. প্রসঙ্গ ও প্রেক্ষাপট (\lat{Context})}

আয়াত ৩৭ সূরার \textbf{\lat{etiquette} (আদব)} অংশের অন্তর্ভুক্ত — যেখানে ইসলামের নৈতিক শিক্ষাগুলো একে একে এসেছে: পিতামাতার সেবা, আত্মীয়ের হক, চলনে শালীনতা। এই আয়াতে আল্লাহ মানুষের দম্ভকে ব্যঙ্গ করেছেন — তুমি অহংকার করে হাঁটো, কিন্তু তুমি তো পৃথিবীকে বিদীর্ণ করতে পারো না, পাহাড়ের উচ্চতায়ও পৌঁছাতে পারো না। অর্থাৎ তোমার শক্তি ও \textbf{\lat{position} (অবস্থান)} সীমিত; তাহলে কেন এত দম্ভ?

\subsubsection*{৩. মূল আরবি}
\begin{center}\arab{وَلَا تَمْشِ فِي الْأَرْضِ مَرَحًا ۖ إِنَّكَ لَن تَخْرِقَ الْأَرْضَ وَلَن تَبْلُغَ الْجِبَالَ طُولًا}\end{center}
\textbf{উচ্চারণ:} ওয়া লা- তামশি ফিল আরদ্বি মারাহা-, ইন্নাকা লান তাখরিক্বাল আরদ্বা ওয়া লান তাবলুগাল জিবা-লা তূলা-।

\subsubsection*{৪. বাংলা অর্থ (\lat{pure})}
\textbf{\lat{“}আর তুমি পৃথিবীতে দম্ভভরে চলো না। তুমি তো কখনো পৃথিবীকে বিদীর্ণ করতে পারবে না, আর উচ্চতায় তুমি পাহাড়কেও ছুঁতে পারবে না।\lat{”}}

\subsubsection*{৫. শব্দের ব্যাখ্যা (\lat{Vocabulary})}

\begin{table}[h]\centering\small
\begin{tabular}{ll}
শব্দ & অর্থ \\
\hline
\textbf{\arab{مَرَحًا}} & দম্ভভরে হাঁটা — স্বাচ্ছন্দ্যের অহংকার \\
\textbf{\arab{تَخْرِقَ}} & বিদীর্ণ করা — ভেদ করা \\
\textbf{\arab{تَبْلُغَ}} & পৌঁছানো \\
\textbf{\arab{طُولًا}} & উচ্চতা — দৈর্ঘ্যে \\
\hline
\end{tabular}
\end{table}

\subsubsection*{৬. গভীর তাফসীর (\lat{Deep} \lat{Tafsir})}

\begin{itemize}
\item \textbf{ইবনে কাসীর:} দম্ভে হাঁটা নিষিদ্ধ করে আল্লাহ বলছেন — তোমার দম্ভ কি পৃথিবী ফুঁড়ে যাবে বা পাহাড়ের মতো উঁচু হবে? বরং তুমি তো মাটির নিচে কবরে মিশে যাবে। এখানে আল্লাহ অহংকারীকে তার \textbf{\lat{end} (পরিণতি)} স্মরণ করিয়ে দিচ্ছেন। (ইবনে কাসীর, ৩/৪৭-৪৮)
\item \textbf{তাবারী:} মানুষকে মনে করিয়ে দেওয়া হচ্ছে তার দুর্বলতা — তার শক্তি দিয়ে সে পৃথিবী ভেদ করতে বা পাহাড়ের উচ্চতায় পৌঁছাতে পারে না, অথচ সে অহংকার করে। (তাবারী, ১৫/৬০)
\item \textbf{কুরতুবী:} আয়াতটি বুকে ভর দিয়ে গর্বভরে হাঁটা নিষেধ করছে। এ ধরনের চলা অহংকারের লক্ষণ, যা আল্লাহ অপছন্দ করেন। (কুরতুবী, ১০/২৮৬)
\item \textbf{মাআরিফুল কুরআন:} এখানে অহংকারকে উপহাস করা হয়েছে — মানুষ তার সীমা ভুলে দম্ভ করে; অথচ সে তো চিরন্তন কিছু নয়, মৃত্যুর পর মাটিতে মিশে যাবে। (মাআরিফুল কুরআন, ৬/৮২)
\end{itemize}
\subsubsection*{৭. সংশ্লিষ্ট হাদিস ও আয়াত}

\begin{itemize}
\item \textbf{আয়াত (লুকমান ৩১:১৮):} \lat{“}আর পৃথিবীতে দম্ভভরে চলো না...\lat{”} — একই শিক্ষা।
\item \textbf{হাদিস (বুখারি ৫৭৮৮):} \lat{“}এক ব্যক্তি অহংকারে হাঁটছিল, আল্লাহ পৃথিবীকে তাকে গ্রাস করার আদেশ দিলেন।\lat{”}
\item \textbf{আয়াত (ফাতির ৩৫:১০):} \lat{“}যে সম্মান (ইজ্জত) চায়, সে জেনে রাখুক — সম্মান সম্পূর্ণ আল্লাহরই।\lat{”}
\end{itemize}
\subsubsection*{৮. গভীর শিক্ষা (\lat{Deep} \lat{Lessons})}

\begin{enumerate}
\item \textbf{মানুষের সীমাবদ্ধতা — দম্ভের জবাব।} তুমি পৃথিবী ভেদ করতে পারো না, পাহাড়কে ছুঁতে পারো না — এতেই তোমার শক্তির \textbf{\lat{limit} (সীমা)} স্পষ্ট। তাহলে অহংকার কিসের?
\item \textbf{আল্লাহর ব্যঙ্গ — শিক্ষামূলক।} আয়াতে অহংকারীকে হাস্যকরভাবে দেখানো হয়েছে — তার দম্ভ তার শক্তির চেয়ে অনেক বড়। এ আল্লাহর \textbf{\lat{teaching} \lat{method} (শিক্ষার পদ্ধতি)}।
\item \textbf{কবরই আসল ঠিকানা।} যতই দম্ভ করো, একদিন মাটিতে মিশে যাবে — এই সত্য স্মরণ অহংকার ধ্বংস করে।
\item \textbf{বিনয় — প্রকৃত শক্তির চিহ্ন।} যার ভেতরে আত্মবিশ্বাস আছে, তার বাহ্যিক দম্ভের প্রয়োজন নেই।
\end{enumerate}
\subsubsection*{৯. জীবনে প্রয়োগ (\lat{Practical} \lat{Application})}

\begin{enumerate}
\item নিজের \textbf{\lat{strength} (শক্তি)} নিয়ে গর্ব করার আগে নিজের দুর্বলতাগুলো স্মরণ করুন।
\item দম্ভ আসলে কবর ও মৃত্যুর কথা \textbf{\lat{remind} (স্মরণ করিয়ে দিন)} নিজেকে।
\item অহংকারী মানুষের সাথে থাকলে তার দম্ভে উত্তেজিত হবেন না — বিনয়ে তার \textbf{\lat{counter} (প্রতিবাদ)} করুন।
\end{enumerate}
\medskip\hrule\medskip

\subsection*{আয়াত ৬ — সূরা আল-আ'রাফ (৭) — আয়াত ৪০}

\subsubsection*{১. সূত্র কার্ড}

\begin{table}[h]\centering\small
\begin{tabular}{ll}
বিষয় & বিবরণ \\
\hline
\textbf{সূরা} & আল-আ'রাফ (৭) — মক্কী \\
\textbf{আয়াত} & ৪০ \\
\textbf{পারা} & ৮ \\
\textbf{মূল বিষয়} & অহংকারী জান্নাতে প্রবেশ করবে না (উট-সূঁচের মিসাল) \\
\hline
\end{tabular}
\end{table}

\subsubsection*{২. প্রসঙ্গ ও প্রেক্ষাপট (\lat{Context})}

সূরার প্রারম্ভে আল্লাহ মানুষের প্রতি দয়া ও সতর্কতা — দুটিই পাঠিয়েছেন। আয়াত ৪০-এ কাফিরদের দুইটি দোষ উল্লেখ করা হয়েছে: (১) আয়াতকে মিথ্যা বলা, (২) তা থেকে অহংকার করা। এই অহংকারের শাস্তি — আকাশের দরজা তাদের জন্য খোলা হবে না, জান্নাতে প্রবেশ অসম্ভব হবে, যতক্ষণ না উট সূঁচের ছিদ্রে ঢোকে। এটি একটি \textbf{\lat{impossible} (অসম্ভব)} ঘটনা — অর্থাৎ তাদের জান্নাত প্রায় অসম্ভব।

\subsubsection*{৩. মূল আরবি}
\begin{center}\arab{إِنَّ الَّذِينَ كَذَّبُوا بِآيَاتِنَا وَاسْتَكْبَرُوا عَنْهَا لَا تُفَتَّحُ لَهُمْ أَبْوَابُ السَّمَاءِ وَلَا يَدْخُلُونَ الْجَنَّةَ حَتَّى يَلِجَ الْجَمَلُ فِي سَمِّ الْخِيَاطِ ۚ وَكَذَلِكَ نَجْزِي الْمُجْرِمِينَ}\end{center}
\textbf{উচ্চারণ:} ইন্নাল্লাযীনা কাযযাবূ বিআ-য়া-তিনা- ওয়াসতাকবারূ 'আনহা- লা- তুফাত্তাহু লাহুম আবওয়া-বুস সামা-য়ি ওয়া লা- ইয়াদখুলূনাল জান্নাতা হাত্তা- ইয়ালিজাল জামালু ফী সাম্মিল খিয়া-ত্ব, ওয়া কাযা-লিকা নাজযিল মুজরিমীন।

\subsubsection*{৪. বাংলা অর্থ (\lat{pure})}
\textbf{\lat{“}নিশ্চয়ই যারা আমার আয়াতকে মিথ্যা বলেছে এবং তা থেকে অহংকার করেছে, তাদের জন্য আকাশের দরজা খোলা হবে না; এবং উট সূঁচের ছিদ্রে ঢোকার আগ পর্যন্ত তারা জান্নাতে প্রবেশ করবে না। এভাবেই আমি অপরাধীদের প্রতিফল দিই।\lat{”}}

\subsubsection*{৫. শব্দের ব্যাখ্যা (\lat{Vocabulary})}

\begin{table}[h]\centering\small
\begin{tabular}{ll}
শব্দ & অর্থ \\
\hline
\textbf{\arab{كَذَّبُوا}} & মিথ্যা বলেছে — মিথ্যারোপ করেছে \\
\textbf{\arab{وَاسْتَكْبَرُوا} \arab{عَنْهَا}} & এবং তা থেকে অহংকার করেছে — \lat{face}-ঘুরিয়েছে \\
\textbf{\arab{يَلِجَ}} & প্রবেশ করা — ঢোকা \\
\textbf{\arab{الْجَمَلُ}} & উট \\
\textbf{\arab{سَمِّ} \arab{الْخِيَاطِ}} & সূঁচের ছিদ্র \\
\hline
\end{tabular}
\end{table}

\subsubsection*{৬. গভীর তাফসীর (\lat{Deep} \lat{Tafsir})}

\begin{itemize}
\item \textbf{ইবনে কাসীর:} তাদের আমল, দোয়া ও আত্মা — কোনো কিছুই আকাশে ওঠবে না, জান্নাতের দরজা তাদের জন্য বন্ধ থাকবে। উট সূঁচের ছিদ্রে ঢোকা যেমন অসম্ভব, তেমনি অহংকারীর জান্নাতে প্রবেশও অসম্ভব — এটা বুঝানোর জন্য \textbf{\lat{example} (উদাহরণ)} দেওয়া হয়েছে। (ইবনে কাসীর, ২/২২৮-২২৯)
\item \textbf{কুরতুবী:} আয়াতের শাস্তি শুধু মুশরিকদের জন্য নয়; সাধারণভাবে যারা হকের সামনে অহংকার করে, তাদের জন্যও প্রযোজ্য। (কুরতুবী, ৭/২৮৩)
\item \textbf{মাআরিফুল কুরআন:} আয়াতে কিবরের শাস্তি স্পষ্ট — অহংকারী জান্নাতে প্রবেশই করতে পারবে না, যতক্ষণ পর্যন্ত ঘটতে-অসম্ভব ঘটনা (উটের সূঁচে ঢোকা) না ঘটে। (মাআরিফুল কুরআন, ৪/৩৫৩)
\end{itemize}
\subsubsection*{৭. সংশ্লিষ্ট হাদিস ও আয়াত}

\begin{itemize}
\item \textbf{হাদিস (মুসলিম ৯১):} \lat{“}যার হৃদয়ে সরিষার দানার কিবর থাকবে, সে জান্নাতে প্রবেশ করবে না।\lat{”} — আয়াতের সাথে হাদিসের সরাসরি \textbf{\lat{connection} (সংযোগ)}।
\item \textbf{আয়াত (যুমার ৩৯:৭২):} অহংকারীদের আবাস জাহান্নাম।
\item \textbf{হাদিস (মুসলিম ২৮৬৫):} \lat{“}জাহান্নাম বলল: আমার মধ্যে অহংকারীরা আসবে।\lat{”}
\end{itemize}
\subsubsection*{৮. গভীর শিক্ষা (\lat{Deep} \lat{Lessons})}

\begin{enumerate}
\item \textbf{মিথ্যা বলার সাথে অহংকারের সম্পর্ক।} আয়াতে দুটি জিনিস একসাথে বলা হয়েছে — আয়াতকে মিথ্যা বলা আর তা থেকে অহংকার করা। অহংকার মানুষকে সত্য \textbf{\lat{reject} (প্রত্যাখ্যান)} করায়।
\item \textbf{আকাশের দরজা বন্ধ — আমল কবুল না হওয়া।} অহংকারীর আমল-দোয়া এমনকি উপরে পৌঁছায় না; তার ভালো কাজও কবুল হয় না।
\item \textbf{অসম্ভবের উদাহরণ — শিক্ষার শক্তি।} \lat{“}উট সূঁচের ছিদ্রে\lat{”} — এই ছবিটি মন থেকে যায়; এটাই কুরআনের \textbf{\lat{unique} (স্বতন্ত্র)} বর্ণনাশৈলী।
\item \textbf{অহংকার — জান্নাতের পথে বাধা।} ঈমান থাকলেও অহংকার থাকলে জান্নাতের দরজা বন্ধ — এখানে নবীজির হাদিসের সতর্কতা প্রযোজ্য।
\end{enumerate}
\subsubsection*{৯. জীবনে প্রয়োগ (\lat{Practical} \lat{Application})}

\begin{enumerate}
\item নিজের জানা সত্যকে মেনে নিতে অস্বীকার করবেন না — এটাই আয়াতের প্রধান শিক্ষা।
\item ভালো আমল করেও নিজের আমল নিয়ে \textbf{\lat{pride} (গর্ব)} করবেন না — আমল কবুল আল্লাহর দয়া।
\item অসম্ভবের উদাহরণ স্মরণে রাখুন — অহংকারের শাস্তি যেন ভুলে না যাই।
\end{enumerate}
\medskip\hrule\medskip

\subsection*{আয়াত ৭ — সূরা আস-সাফফাত (৩৭) — আয়াত ৩৫}

\subsubsection*{১. সূত্র কার্ড}

\begin{table}[h]\centering\small
\begin{tabular}{ll}
বিষয় & বিবরণ \\
\hline
\textbf{সূরা} & আস-সাফফাত (৩৭) — মক্কী \\
\textbf{আয়াত} & ৩৫ \\
\textbf{পারা} & ২৩ \\
\textbf{মূল বিষয়} & কালেমা শুনে অহংকার — কিবরের চূড়া \\
\hline
\end{tabular}
\end{table}

\subsubsection*{২. প্রসঙ্গ ও প্রেক্ষাপট (\lat{Context})}

আয়াত ৩৩-৩৫-এ কাফিরদের পরিণতি ও তাদের অপরাধের বর্ণনা। আয়াত ৩৫-এ স্পষ্ট — কুরাইশদের সবচেয়ে বড় অপরাধ এই যে, যখন তাদের বলা হতো \lat{“}\arab{لا} \arab{إله} \arab{إلا} \arab{الله}\lat{”} (আল্লাহ ছাড়া কোনো উপাস্য নেই), তখন তারা অহংকার করত। অর্থাৎ সবচেয়ে বড় হক — তাওহীদ — মেনে নিতে তাদের অহংকার বাধা দিয়েছিল।

\subsubsection*{৩. মূল আরবি}
\begin{center}\arab{إِنَّهُمْ كَانُوا إِذَا قِيلَ لَهُمْ لَا إِلَٰهَ إِلَّا اللَّهُ يَسْتَكْبِرُونَ}\end{center}
\textbf{উচ্চারণ:} ইন্নাহুম কা-নূ ইযা- ক্বীলা লাহুম লা- ইলা-হা ইল্লাল্লা-হু ইয়াসতাকবিরূন।

\subsubsection*{৪. বাংলা অর্থ (\lat{pure})}
\textbf{\lat{“}তারা এমন ছিল যে, যখন তাদের বলা হতো — ‘আল্লাহ ছাড়া কোনো উপাস্য নেই’, তখন তারা অহংকার করত।\lat{”}}

\subsubsection*{৫. শব্দের ব্যাখ্যা (\lat{Vocabulary})}

\begin{table}[h]\centering\small
\begin{tabular}{ll}
শব্দ & অর্থ \\
\hline
\textbf{\arab{لَا} \arab{إِلَٰهَ} \arab{إِلَّا} \arab{اللَّهُ}} & আল্লাহ ছাড়া কোনো উপাস্য নেই — কালেমা \\
\textbf{\arab{يَسْتَكْبِرُونَ}} & অহংকার করত — মাথা নত করতে অস্বীকার \\
\hline
\end{tabular}
\end{table}

\subsubsection*{৬. গভীর তাফসীর (\lat{Deep} \lat{Tafsir})}

\begin{itemize}
\item \textbf{ইবনে কাসীর:} কুরাইশরা তাওহীদের কালিমা শুনলে অহংকারে তা মেনে নিত না — এটাই তাদের কিবর। কিবরের কারণে তারা জান্নাত থেকে বঞ্চিত হয়েছিল। (ইবনে কাসীর, ৪/১৭-১৮)
\item \textbf{বাগাবী:} কালেমা শুনে অহংকার করা — এটি কিবরের সবচেয়ে বড় প্রকার; কারণ নিজেকে বড় মনে করে আল্লাহর বড়ত্ব মানতে চায় না। (বাগাবী, ৪/১৮৮)
\item \textbf{মাআরিফুল কুরআন:} \lat{“}\arab{لا} \arab{إله} \arab{إلا} \arab{الله}\lat{”} শুনে অহংকার — এ থেকে বুঝা যায় কিবর কুফর ও শিরকের মূল উপাদান। আল্লাহর হুকুমের সামনে নিজের গর্বকে নত করতে পারা বিনয়ের প্রথম শর্ত। (মাআরিফুল কুরআন, ৬/৪০৬)
\end{itemize}
\subsubsection*{৭. সংশ্লিষ্ট হাদিস ও আয়াত}

\begin{itemize}
\item \textbf{আয়াত (গাফির ৪০:৬০):} \lat{“}যারা আমার ইবাদত থেকে অহংকার করে, তারা লাঞ্ছিত হয়ে জাহান্নামে প্রবেশ করবে।\lat{”}
\item \textbf{হাদিস (মুসলিম ৯১):} \lat{“}কিবর হলো সত্যকে তুচ্ছ করা।\lat{”} — কালেমা মেনে না নেওয়া সত্য তুচ্ছ করার চরম রূপ।
\item \textbf{আয়াত (আ'রাফ ৭:১৪৬):} অহংকারীদের থেকে সত্যবোধ কেড়ে নেওয়া।
\end{itemize}
\subsubsection*{৮. গভীর শিক্ষা (\lat{Deep} \lat{Lessons})}

\begin{enumerate}
\item \textbf{সবচেয়ে বড় হকের সামনে অহংকার — চরম কিবর।} মানুষ সামাজিক, আর্থিক বা জ্ঞানগত বিষয়ে অহংকার করে; কিন্তু সবচেয়ে মারাত্মক — আল্লাহর সত্যের সামনে অহংকার।
\item \textbf{অহংকার — ঈমানের প্রতিবন্ধক।} কুরাইশরা ঈমান আনে না কারণ তাদের অহংকার কালেমা মেনে নিতে দেয়নি। অহংকার ও ঈমান একসাথে থাকে না।
\item \textbf{বিনয় — ঈমানের প্রথম ধাপ।} ঈমান মানেই আল্লাহর সামনে নিজেকে নত করা — এটাই বিনয়ের শুরু।
\item \textbf{কিবর কুফর-শিরকের মূল।} শিরকও মূলত অহংকারেরই রূপ — নিজের \textbf{\lat{opinion} (মত)} আল্লাহর আদেশের উপরে রাখা।
\end{enumerate}
\subsubsection*{৯. জীবনে প্রয়োগ (\lat{Practical} \lat{Application})}

\begin{enumerate}
\item নিজেকে প্রশ্ন করুন: \lat{“}আমি কি আল্লাহর কোনো আদেশ মেনে নিতে অহংকার করি?\lat{”}
\item তাওহীদের শিক্ষা শিখে অন্যকে \textbf{\lat{share} (ভাগ করে নিন)} করুন — জ্ঞানে বিনয়ী হোন।
\item সত্য মেনে নিতে দেরি করলে তাড়াতাড়ি করুন — অহংকার যত বাড়ে, মেনে নেওয়া তত কঠিন।
\end{enumerate}
\medskip\hrule\medskip

\subsection*{আয়াত ৮ — সূরা আল-বাকারাহ (২) — আয়াত ২০৬}

\subsubsection*{১. সূত্র কার্ড}

\begin{table}[h]\centering\small
\begin{tabular}{ll}
বিষয় & বিবরণ \\
\hline
\textbf{সূরা} & আল-বাকারাহ (২) — মাদানী \\
\textbf{আয়াত} & ২০৬ \\
\textbf{পারা} & ২ \\
\textbf{মূল বিষয়} & নসীহত শুনে অহংকার — কিবরের লক্ষণ \\
\hline
\end{tabular}
\end{table}

\subsubsection*{২. প্রসঙ্গ ও প্রেক্ষাপট (\lat{Context})}

আয়াত ২০৪-২০৭ মুনাফিক ও মুমিনের পার্থক্য বর্ণনা করছে। আয়াত ২০৪-এ মুনাফিকের চেহারা — দুনিয়ার কথায় মানুষকে মুগ্ধ করে, কিন্তু ভেতরে শত্রুতা। আয়াত ২০৫-এ তার ধ্বংসাত্মক কাজ। আয়াত ২০৬-এ তার চরম দোষ — যখন তাকে \textbf{নসীহত (উপদেশ)} করা হয় যে \lat{“}আল্লাহকে ভয় করো\lat{”}, তখন তার অহংকার তাকে পাপের দিকে টেনে নেয়। অর্থাৎ নসীহত মেনে নিতে তার অনীহাই অহংকার।

\subsubsection*{৩. মূল আরবি}
\begin{center}\arab{وَإِذَا قِيلَ لَهُ اتَّقِ اللَّهَ أَخَذَتْهُ الْعِزَّةُ بِالْإِثْمِ ۚ فَحَسْبُهُ جَهَنَّمُ ۚ وَلَبِئْسَ الْمِهَادُ}\end{center}
\textbf{উচ্চারণ:} ওয়া ইযা- ক্বীলা লাহুত্তাক্বিল্লা-হা আখাযাতহুল 'ইযযাতু বিল ইছমি, ফাহাসবুহূ জাহান্নামু ওয়া লাবি'সাল মিহা-দ।

\subsubsection*{৪. বাংলা অর্থ (\lat{pure})}
\textbf{\lat{“}আর যখন তাকে বলা হয় — ‘আল্লাহকে ভয় করো’, তখন গর্ব তাকে পাপের দিকে টেনে নেয় (অর্থাৎ অহংকারবশে সে নসীহত প্রত্যাখ্যান করে)। তার জন্য যথেষ্ট জাহান্নাম; আর কতই না নিকৃষ্ট সেই বিশ্রামস্থল!\lat{”}}

\subsubsection*{৫. শব্দের ব্যাখ্যা (\lat{Vocabulary})}

\begin{table}[h]\centering\small
\begin{tabular}{ll}
শব্দ & অর্থ \\
\hline
\textbf{\arab{اتَّقِ} \arab{اللَّهَ}} & আল্লাহকে ভয় করো — তাকওয়া অবলম্বন করো \\
\textbf{\arab{أَخَذَتْهُ} \arab{الْعِزَّةُ} \arab{بِالْإِثْمِ}} & গর্ব তাকে পাপে টেনে নেয় — অহংকার \\
\textbf{\arab{حَسْبُهُ} \arab{جَهَنَّمُ}} & জাহান্নামই তার জন্য যথেষ্ট \\
\textbf{\arab{الْمِهَادُ}} & বিশ্রামস্থল — শয্যা \\
\hline
\end{tabular}
\end{table}

\subsubsection*{৬. গভীর তাফসীর (\lat{Deep} \lat{Tafsir})}

\begin{itemize}
\item \textbf{ইবনে কাসীর:} কেউ কাউকে আল্লাহর তাকওয়া অবলম্বনে নসীহত করলে সে রাগ ও অহংকারে তা প্রত্যাখ্যান করে — এটা মুনাফিকের চরিত্র। (ইবনে কাসীর, ১/৩৪৬)
\item \textbf{তাবারী:} \lat{“}\arab{أَخَذَتْهُ} \arab{الْعِزَّةُ} \arab{بِالْإِثْمِ}\lat{”} অর্থ — নসীহত শুনে তার অহংকার জেগে ওঠে, ফলে সে পাপ ও অন্যায়ের পথেই থেকে যায়। (তাবারী, ২/৩১০)
\item \textbf{সা'দী:} নসীহতকারীকে ছোট করে দেখা, নসীহতকে মানতে অস্বীকৃতি — এটাই কিবর। এর প্রতিদান জাহান্নাম। (আস-সা'দী, পৃষ্ঠা ৯১)
\item \textbf{মাআরিফুল কুরআন:} এই আয়াত স্পষ্ট করে — সৎ উপদেশ শুনে অহংকার করা কিবর, যা মানুষকে জাহান্নামে নিয়ে যায়। কিবরের প্রথম লক্ষণ নসীহত গ্রহণে অনীহা। (মাআরিফুল কুরআন, ১/২৭৮)
\end{itemize}
\subsubsection*{৭. সংশ্লিষ্ট হাদিস ও আয়াত}

\begin{itemize}
\item \textbf{হাদিস (আবু দাউদ ৫১২৪):} \lat{“}দ্বীন হলো নসীহত।\lat{”} — নসীহত গ্রহণ করা দ্বীনের অংশ।
\item \textbf{আয়াত (যুমার ৩৯:৭২):} অহংকারীদের আবাস জাহান্নাম।
\item \textbf{হাদিস (মুসলিম ৯১):} \lat{“}কিবর হলো সত্যকে তুচ্ছ করা।\lat{”} — নসীহত মেনে না নেওয়া সত্য তুচ্ছ করা।
\end{itemize}
\subsubsection*{৮. গভীর শিক্ষা (\lat{Deep} \lat{Lessons})}

\begin{enumerate}
\item \textbf{নসীহত গ্রহণের অনীহা — কিবরের প্রথম লক্ষণ।} মানুষ ধর্মীয় নসীহত, বড়দের উপদেশ, বন্ধুর পরামর্শ — মেনে নিতে অহংকার করে। এটাই আয়াতের মূল বার্তা।
\item \textbf{অহংকার মানুষকে পাপের দিকে টানে।} \lat{“}আল্লাহকে ভয় করো\lat{”} শুনে রাগ না করে মেনে নিলে পাপ বন্ধ হতো; অহংকার তা দেরি করায়।
\item \textbf{জাহান্নামই যথেষ্ট।} অহংকারীকে কোনো দীর্ঘ শাস্তির তালিকা দেওয়া হয়নি; বলা হয়েছে — জাহান্নামই তার জন্য \textbf{\lat{enough} (যথেষ্ট)}।
\item \textbf{মুমিনের চিহ্ন — উপদেশ মেনে নেওয়া।} মুনাফিকের বিপরীত চরিত্র হলো মুমিনের — যিনি সত্য শুনে বিনয়ে মেনে নেন।
\end{enumerate}
\subsubsection*{৯. জীবনে প্রয়োগ (\lat{Practical} \lat{Application})}

\begin{enumerate}
\item কোনো উপদেশ শুনে প্রথম \textbf{\lat{reaction} (প্রতিক্রিয়া)} নজর করুন — রাগ আসছে না বিনয় আসছে?
\item ছোটদের কাছ থেকে উপদেশ পেলেও \textbf{\lat{listen} (শুনুন)} — হক ছোট-বড় সবার কাছ থেকে আসতে পারে।
\item প্রতিদিন একজন বন্ধুকে সত্য কথা বলুন ও সত্য কথা শুনুন — এতে অহংকার ভাঙবে।
\end{enumerate}
\subsection*{আয়াত ৯ — সূরা আল-আ'রাফ (৭) — আয়াত ১৪৬}

\subsubsection*{১. সূত্র কার্ড}

\begin{table}[h]\centering\small
\begin{tabular}{ll}
বিষয় & বিবরণ \\
\hline
\textbf{সূরা} & আল-আ'রাফ (৭) — মক্কী \\
\textbf{আয়াত} & ১৪৬ \\
\textbf{পারা} & ৯ \\
\textbf{মূল বিষয়} & অহংকারীর অন্তর থেকে সত্যবোধ কেড়ে নেওয়া হয় \\
\hline
\end{tabular}
\end{table}

\subsubsection*{২. প্রসঙ্গ ও প্রেক্ষাপট (\lat{Context})}

আয়াত ১৪৫-এ আল্লাহ মূসা (আঃ)-কে তাওরাতের ফলক ও তার বিধান দেওয়ার কথা বলেছেন। এরপর আয়াত ১৪৬-এ আল্লাহ বলছেন — যারা পৃথিবীতে অহংকার করে, আমি তাদেরকে আমার নিদর্শন থেকে দূরে সরিয়ে দেবো। এটাই অহংকারের সবচেয়ে বড় শাস্তি — আল্লাহর হেদায়াত থেকে বঞ্চিত হওয়া।

\subsubsection*{৩. মূল আরবি}
\begin{center}\arab{سَأَصْرِفُ عَنْ آيَاتِيَ الَّذِينَ يَتَكَبَّرُونَ فِي الْأَرْضِ بِغَيْرِ الْحَقِّ ۚ وَإِن يَرَوْا كُلَّ آيَةٍ لَّا يُؤْمِنُوا بِهَا ۚ وَإِن يَرَوْا سَبِيلَ الرُّشْدِ لَا يَتَّخِذُوهُ سَبِيلًا ۚ وَإِن يَرَوْا سَبِيلَ الْغَيِّ يَتَّخِذُوهُ سَبِيلًا}\end{center}
\textbf{উচ্চারণ:} সাআসরিফু 'আন আ-য়া-তিয়াল্লাযীনা ইয়াতাকাব্বারূনা ফিল আরদ্বি বিগাইরিল হাক্বক্বি, ওয়া ইন ইয়ারাও কুল্লা আ-য়াতিল লা- ইউ'মিনূ বিহা-, ওয়া ইন ইয়ারাও সাবীলার রুশদি লা- ইয়াত্তাখিযূহু সাবীলা-, ওয়া ইন ইয়ারাও সাবীলাল গায়্যি ইয়াত্তাখিযূহু সাবীলা-।

\subsubsection*{৪. বাংলা অর্থ (\lat{pure})}
\textbf{\lat{“}অচিরেই আমি আমার নিদর্শনগুলো থেকে দূরে সরিয়ে দেবো তাদেরকে যারা পৃথিবীতে অন্যায়ভাবে অহংকার করে। আর তারা যদি প্রতিটি নিদর্শনও দেখে, তবুও তাতে ঈমান আনবে না; যদি তারা সঠিক পথ দেখে, তাকে নিজের পথ বানাবে না; আর যদি ভ্রান্ত পথ দেখে, সেটাকেই নিজের পথ বানাবে।\lat{”}}

\subsubsection*{৫. শব্দের ব্যাখ্যা (\lat{Vocabulary})}

\begin{table}[h]\centering\small
\begin{tabular}{ll}
শব্দ & অর্থ \\
\hline
\textbf{\arab{سَأَصْرِفُ}} & আমি দূরে সরিয়ে দেবো — \lat{deprive} \\
\textbf{\arab{يَتَكَبَّرُونَ}} & অহংকার করে \\
\textbf{\arab{سَبِيلَ} \arab{الرُّشْدِ}} & সঠিক/সরল পথ \\
\textbf{\arab{سَبِيلَ} \arab{الْغَيِّ}} & ভ্রান্ত/বিপথ \\
\textbf{\arab{آيَاتِيَ}} & আমার নিদর্শনসমূহ \\
\hline
\end{tabular}
\end{table}

\subsubsection*{৬. গভীর তাফসীর (\lat{Deep} \lat{Tafsir})}

\begin{itemize}
\item \textbf{ইবনে কাসীর:} অহংকারীদের থেকে আল্লাহ হক বুঝার তাওফিক কেড়ে নেন। এটা কিবরের কঠিন শাস্তি — সত্য দেখেও সত্য গ্রহণ না করা। (ইবনে কাসীর, ২/২৫৩-২৫৪)
\item \textbf{তাবারী:} অহংকারীদের অন্তর থেকে সত্য উপলব্ধির জ্ঞান কেড়ে নেওয়া হবে, ফলে তারা সত্য চেনার পরও সেটা বর্জন করবে। (তাবারী, ৫/৮২)
\item \textbf{মাআরিফুল কুরআন:} অহংকার যাদের বৈশিষ্ট্য, আল্লাহ তাদের অন্তর থেকে হেদায়াতের নূর কেড়ে নেন। ফলে দুনিয়ার ক্ষমতা-পদ তাদের ভোগ করবে কিন্তু আখেরাত বিফল। (মাআরিফুল কুরআন, ৪/৪৩৩)
\end{itemize}
\subsubsection*{৭. সংশ্লিষ্ট হাদিস ও আয়াত}

\begin{itemize}
\item \textbf{আয়াত (বাকারাহ ২:২০৬):} নসীহত শুনে অহংকার — কিবরের লক্ষণ।
\item \textbf{হাদিস (বুখারি ৭৫৯৬):} \lat{“}যে জ্ঞান কেবল মানুষের কাছে বড় দেখানোর জন্য অর্জন করে, সে জান্নাতের ঘ্রাণও পাবে না।\lat{”} — নিদর্শন থেকে বঞ্চিত হওয়ার কাছাকাছি সতর্কতা।
\item \textbf{আয়াত (নাহল ১৬:২৩):} \lat{“}আল্লাহ অহংকারীদের পছন্দ করেন না।\lat{”}
\end{itemize}
\subsubsection*{৮. গভীর শিক্ষা (\lat{Deep} \lat{Lessons})}

\begin{enumerate}
\item \textbf{সবচেয়ে বড় শাস্তি — হেদায়াত থেকে বঞ্চিত হওয়া।} ধন-সম্পদ না হারিয়ে আল্লাহর নিদর্শন বুঝার শক্তি হারানোই সবচেয়ে কঠিন শাস্তি।
\item \textbf{অহংকারের চক্র।} অহংকারী সত্য দেখেও মেনে নেয় না, ভুল পথ দেখলে তাকে পথ বানায় — এটা একটা \textbf{\lat{vicious} \lat{cycle} (দুষ্টচক্র)}।
\item \textbf{বোঝার শক্তি আল্লাহর দান।} যে বুঝার শক্তি আছে, তা আল্লাহর \textbf{\lat{blessing} (নেয়ামত)}; অহংকার তা কেড়ে নিতে পারে।
\item \textbf{সত্য গ্রহণে তাড়াতাড়ি করুন।} অহংকার যত বাড়ে, সত্য গ্রহণ তত কঠিন — তাই দেরি না করে সত্য মেনে নিন।
\end{enumerate}
\subsubsection*{৯. জীবনে প্রয়োগ (\lat{Practical} \lat{Application})}

\begin{enumerate}
\item সত্যের পথ দেখলে সাথে সাথে সেটি গ্রহণ করুন — \textbf{\lat{delay} (দেরি)} অহংকারের লক্ষণ হতে পারে।
\item অহংকারের কারণে সত্য মেনে না নেওয়ার \textbf{\lat{habit} (অভ্যাস)} থাকলে তা ভাঙুন — দোয়া করুন হেদায়াতের জন্য।
\item জ্ঞান ও বুঝের জন্য আল্লাহর কাছে শুকরিয়া \textbf{\lat{express} (প্রকাশ)} করুন।
\end{enumerate}
\medskip\hrule\medskip

\subsection*{আয়াত ১০ — সূরা আল-কাসাস (২৮) — আয়াত ৮৩}

\subsubsection*{১. সূত্র কার্ড}

\begin{table}[h]\centering\small
\begin{tabular}{ll}
বিষয় & বিবরণ \\
\hline
\textbf{সূরা} & আল-কাসাস (২৮) — মক্কী \\
\textbf{আয়াত} & ৮৩ \\
\textbf{পারা} & ২০ \\
\textbf{মূল বিষয়} & আখেরাত তাদের, যারা উচ্চতা ও ফাসাদ চায় না \\
\hline
\end{tabular}
\end{table}

\subsubsection*{২. প্রসঙ্গ ও প্রেক্ষাপট (\lat{Context})}

সূরায় কারুনের ঘটনা এসেছে — ধনী ব্যক্তি যার ধন-সম্পদ ছিল অপার, কিন্তু সে অহংকার করে। আয়াত ৭৬-৮২-এ কারুনের অহংকার, ধন-সম্পদে গর্ব ও শেষে পৃথিবীতে ধ্বসে যাওয়ার বর্ণনা। তারপর আয়াত ৮৩-এ আল্লাহ \textbf{\lat{final} \lat{verdict} (চূড়ান্ত সিদ্ধান্ত)} দিচ্ছেন — আখেরাতের ঘর তাদের জন্য, যারা পৃথিবীতে উচ্চতা (বড়ত্ব) ও ফাসাদ (বিপর্যয়) কামনা করে না। কারুনের গল্পের \textbf{\lat{moral} (নৈতিক শিক্ষা)} — অহংকার ও ফাসাদই ধ্বংসের কারণ।

\subsubsection*{৩. মূল আরবি}
\begin{center}\arab{تِلْكَ الدَّارُ الْآخِرَةُ نَجْعَلُهَا لِلَّذِينَ لَا يُرِيدُونَ عُلُوًّا فِي الْأَرْضِ وَلَا فَسَادًا ۚ وَالْعَاقِبَةُ لِلْمُتَّقِينَ}\end{center}
\textbf{উচ্চারণ:} তিলকাদ দা-রুল আ-খিরাতু নাজ'আলুহা- লিল্লাযীনা লা- ইউরীদূনা 'উলুওওয়ান ফিল আরদ্বি ওয়া লা- ফাসা-দা-, ওয়াল 'আ-ক্বিবাতু লিলমুত্তাক্বীন।

\subsubsection*{৪. বাংলা অর্থ (\lat{pure})}
\textbf{\lat{“}এ সেই আখেরাতের ঘর, যা আমি তাদেরকে দেবো যারা পৃথিবীতে উচ্চতা (বড়ত্ব) ও বিপর্যয় কামনা করে না। আর উত্তম পরিণতি তো মুত্তাকীদের জন্যই।\lat{”}}

\subsubsection*{৫. শব্দের ব্যাখ্যা (\lat{Vocabulary})}

\begin{table}[h]\centering\small
\begin{tabular}{ll}
শব্দ & অর্থ \\
\hline
\textbf{\arab{الدَّارُ} \arab{الْآخِرَةُ}} & আখেরাতের ঘর \\
\textbf{\arab{عُلُوًّا}} & উচ্চতা — বড়ত্ব, \lat{dominance} \\
\textbf{\arab{فَسَادًا}} & বিপর্যয় — অন্যায় ও জুলুম \\
\textbf{\arab{الْعَاقِبَةُ}} & চূড়ান্ত পরিণতি — \lat{final} \lat{outcome} \\
\hline
\end{tabular}
\end{table}

\subsubsection*{৬. গভীর তাফসীর (\lat{Deep} \lat{Tafsir})}

\begin{itemize}
\item \textbf{ইবনে কাসীর:} আখেরাতের সফলতা সেই মুত্তাকীদের জন্য যারা জমিনে গর্ব, স্বেচ্ছাচারিতা ও ফাসাদ কামনা করে না। এটা বিনয়ের পুরস্কার। (ইবনে কাসীর, ৩/৩২৭)
\item \textbf{তাবারী:} \lat{“}\arab{عُلُوًّا}\lat{”} অর্থ বড়ত্ব ও গর্ব, \lat{“}\arab{فَسَادًا}\lat{”} অর্থ অন্যায় ও জুলুম। যে এই দুটি থেকে মুক্ত, তার জন্যই আখেরাতের কল্যাণ। (তাবারী, ২০/৯০)
\item \textbf{মাআরিফুল কুরআন:} আখেরাতের মর্যাদা পাওয়ার শর্ত হিসেবে আল্লাহ বড়ত্ব (উলু) ও বিপর্যয় (ফাসাদ) ত্যাগের কথা বলেছেন — উলু ত্যাগই কিবর ত্যাগের পরীক্ষা। (মাআরিফুল কুরআন, ৭/৩৫০)
\end{itemize}
\subsubsection*{৭. সংশ্লিষ্ট হাদিস ও আয়াত}

\begin{itemize}
\item \textbf{হাদিস (মুসলিম ২৫৮৮):} \lat{“}যে আল্লাহর জন্য বিনয় করে, আল্লাহ তাকে উন্নত করেন।\lat{”} — আয়াতের প্রতিদান।
\item \textbf{হাদিস (তিরমিজি ২০২৯):} \lat{“}যে নিজেকে আল্লাহর জন্য ছোট করে, আল্লাহ তাকে সম্মানিত করেন।\lat{”}
\item \textbf{আয়াত (যুমার ৩৯:৭২):} অহংকারীদের আবাস জাহান্নাম — বিপরীত দিকের সতর্কতা।
\end{itemize}
\subsubsection*{৮. গভীর শিক্ষা (\lat{Deep} \lat{Lessons})}

\begin{enumerate}
\item \textbf{বিনয় — আখেরাতের চাবিকাঠি।} যারা দুনিয়ায় বড় হতে চায় না, আল্লাহ তাদেরকে আখেরাতে বড় করেন।
\item \textbf{উলু ও ফাসাদ একসাথে।} অহংকার (উলু) প্রায়ই অন্যায় (ফাসাদ)-এর দিকে নিয়ে যায়; বড় হওয়ার ইচ্ছা মানুষকে জুলুম করতে বাধ্য করে।
\item \textbf{শক্তিশালী প্রতিদান — আখেরাতের ঘর।} ধন-সম্পদ, ক্ষমতা — এগুলো দুনিয়ার; আসল সম্পদ আখেরাতের স্থায়ী ঘর।
\item \textbf{কারুনের বিপরীত চিত্র।} কারুন অহংকারে ধ্বংস; মুত্তাকী বিনয়ে সফল — দুই চিত্র এক আয়াতে।
\end{enumerate}
\subsubsection*{৯. জীবনে প্রয়োগ (\lat{Practical} \lat{Application})}

\begin{enumerate}
\item বড় পদ বা সম্পদ পেলে তার জন্য \textbf{\lat{proud} (গর্ব)} না হয়ে আল্লাহর \textbf{\lat{shukr} (শুকরিয়া)} করুন।
\item অন্যায়ের \textbf{\lat{temptation} (প্রলোভন)} এলে মনে রাখুন — আখেরাতের ঘর তাদের, যারা ফাসাদ চায় না।
\item বিনয়ী মানুষদের সাথে থাকুন — তাদের চরিত্র \textbf{\lat{inspire} (অনুপ্রাণিত)} করে।
\end{enumerate}
\medskip\hrule\medskip

\subsection*{আয়াত ১১ — সূরা আন-নিসা (৪) — আয়াত ৩৬ (শেষাংশ)}

\subsubsection*{১. সূত্র কার্ড}

\begin{table}[h]\centering\small
\begin{tabular}{ll}
বিষয় & বিবরণ \\
\hline
\textbf{সূরা} & আন-নিসা (৪) — মাদানী \\
\textbf{আয়াত} & ৩৬ (শেষাংশ) \\
\textbf{পারা} & ৫ \\
\textbf{মূল বিষয়} & দাম্ভিক-আত্মপ্রশংসাকারী আল্লাহর অপছন্দ \\
\hline
\end{tabular}
\end{table}

\subsubsection*{২. প্রসঙ্গ ও প্রেক্ষাপট (\lat{Context})}

আয়াত ৩৬ ইসলামের সবচেয়ে বিস্তৃত সামাজিক \textbf{\lat{rights} (অধিকার)} আয়াত — এতে আল্লাহর ইবাদত, শিরক থেকে বাঁচা, পিতামাতা, আত্মীয়, ইয়াতিম, মিসকিন, প্রতিবেশী, সফরসঙ্গী — সবার হক আদায়ের নির্দেশ। আয়াতের শেষে আল্লাহ সতর্ক করেছেন — এই হকগুলো আদায়ের সময় যদি অহংকার ও আত্মপ্রশংসা জেগে ওঠে, তবে তা আল্লাহর অপছন্দ। অর্থাৎ ভালো কাজ করার সময়ও অহংকার যেন না আসে।

\subsubsection*{৩. মূল আরবি}
\begin{center}\arab{إِنَّ اللَّهَ لَا يُحِبُّ مَن كَانَ مُخْتَالًا فَخُورًا}\end{center}
\textbf{উচ্চারণ:} ইন্নাল্লা-হা লা- ইউহিব্বু মান কা-না মুখতা-লান ফাখূরা-।

\subsubsection*{৪. বাংলা অর্থ (\lat{pure})}
\textbf{\lat{“}নিশ্চয়ই আল্লাহ তাদেরকে পছন্দ করেন না যারা দাম্ভিক (মুখতাল) ও আত্মপ্রশংসাকারী (ফাখুর)।\lat{”}}

\subsubsection*{৫. শব্দের ব্যাখ্যা (\lat{Vocabulary})}

\begin{table}[h]\centering\small
\begin{tabular}{ll}
শব্দ & অর্থ \\
\hline
\textbf{\arab{مُخْتَالًا}} & দাম্ভিক — নিজেকে বড় ভাবা \\
\textbf{\arab{فَخُورًا}} & আত্মপ্রশংসাকারী — নিজের গুণাগুণ বর্ণনা করে বেড়ানো \\
\textbf{\arab{لَا} \arab{يُحِبُّ}} & ভালোবাসেন না — অপছন্দ \\
\hline
\end{tabular}
\end{table}

\subsubsection*{৬. গভীর তাফসীর (\lat{Deep} \lat{Tafsir})}

\begin{itemize}
\item \textbf{ইবনে কাসীর:} \lat{“}\arab{مُخْتَال}\lat{”} — যে অহংকারে চলাফেরা করে; \lat{“}\arab{فَخُور}\lat{”} — যে নিজের প্রশংসা করে ও মানুষকে হেয় জ্ঞান করে। আয়াতে আল্লাহ এই দুটি গুণের লোককে অপছন্দের কথা বলেছেন। (ইবনে কাসীর, ১/৪৪৬)
\item \textbf{সা'দী:} মুখতাল-ফাখুর হলো কিবর ও উজবের সংমিশ্রণ — নিজেকে বড় দেখানো আর নিজের প্রশংসা করা। এই আয়াতে আত্মীয়তার হক, ইয়াতিম-মিসকিনের হক আদায়ের পর এদের থেকেই সাবধান করা হচ্ছে। (আস-সা'দী, পৃষ্ঠা ১৭৫)
\item \textbf{মাআরিফুল কুরআন:} আয়াতের প্রেক্ষাপট — যখন আত্মীয়-প্রতিবেশীর হক আদায় করা হয়, তখনই প্রবলভাবে কিবর ও উজব জেগে ওঠার আশঙ্কা; তাই আল্লাহ নিজেই বলছেন — দাম্ভিক-আত্মপ্রশংসাকারীকে তিনি পছন্দ করেন না। (মাআরিফুল কুরআন, ২/১২০)
\end{itemize}
\subsubsection*{৭. সংশ্লিষ্ট হাদিস ও আয়াত}

\begin{itemize}
\item \textbf{আয়াত (হাদীদ ৫৭:২৩):} \lat{“}আর আল্লাহ কোনো দাম্ভিক, আত্মপ্রশংসাকারীকে পছন্দ করেন না।\lat{”} — একই বাক্য।
\item \textbf{হাদিস (মুসলিম ৯১):} \lat{“}যার হৃদয়ে সরিষার দানার কিবর থাকবে, সে জান্নাতে প্রবেশ করবে না।\lat{”}
\item \textbf{আয়াত (লুকমান ৩১:১৮):} \lat{“}নিশ্চয়ই আল্লাহ কোনো আত্মগর্বী, দাম্ভিক ব্যক্তিকে পছন্দ করেন না।\lat{”}
\end{itemize}
\subsubsection*{৮. গভীর শিক্ষা (\lat{Deep} \lat{Lessons})}

\begin{enumerate}
\item \textbf{ভালো কাজেও অহংকারের আশঙ্কা।} আয়াতটি সামাজিক হক আদায়ের প্রসঙ্গে এসেছে — অর্থাৎ মানুষ ভালো কাজ করেও অহংকার ও আত্মপ্রশংসায় পড়তে পারে।
\item \textbf{মুখতাল ও ফাখুর — দুই রকমের অহংকার।} একজন নিজের \textbf{\lat{position} (অবস্থান)} নিয়ে গর্ব করে; অন্যজন নিজের গুণ নিয়ে কথা বলে বেড়ায়।
\item \textbf{হক আদায় — নিয়মিত আমল।} পিতামাতা, আত্মীয়, প্রতিবেশী — এদের হক আদায় অহংকার নয়, বিনয়ের শিক্ষা।
\item \textbf{অপছন্দ থেকে বাঁচার পথ — ইখলাস।} যার আমল শুধু আল্লাহর জন্য, তার প্রশংসা মানুষের কাছে চাওয়ার দরকার নেই।
\end{enumerate}
\subsubsection*{৯. জীবনে প্রয়োগ (\lat{Practical} \lat{Application})}

\begin{enumerate}
\item দান-সদকা বা ভালো কাজের পর নিজের প্রশংসা \textbf{\lat{mention} (উল্লেখ)} করবেন না — এতে আমল নষ্ট হয়।
\item নিজের \textbf{\lat{achievement} (অর্জন)} নিয়ে কথা বলার আগে ভাবুন — এটা কি আল্লাহর দেখানোর জন্য নাকি মানুষের কাছে বড় দেখানোর জন্য?
\item প্রতিবেশী, আত্মীয়, ইয়াতিম-মিসকিনের হক আদায়ে \textbf{\lat{regular} (নিয়মিত)} থাকুন — এতে উজব কমে।
\end{enumerate}
\medskip\hrule\medskip

\subsection*{আয়াত ১২ — সূরা আল-হাদীদ (৫৭) — আয়াত ২৩}

\subsubsection*{১. সূত্র কার্ড}

\begin{table}[h]\centering\small
\begin{tabular}{ll}
বিষয় & বিবরণ \\
\hline
\textbf{সূরা} & আল-হাদীদ (৫৭) — মাদানী \\
\textbf{আয়াত} & ২৩ \\
\textbf{পারা} & ২৭ \\
\textbf{মূল বিষয়} & প্রাপ্তিতে দম্ভ (উজব) নিষিদ্ধ \\
\hline
\end{tabular}
\end{table}

\subsubsection*{২. প্রসঙ্গ ও প্রেক্ষাপট (\lat{Context})}

আয়াত ২২-২৩-এ আল্লাহ বলছেন — পৃথিবীতে যা কিছু ঘটে, সব আল্লাহর কিতাবে লেখা আছে, ঘটার আগেই। তারপর আয়াত ২৩-এ এর ফল — তোমরা যা হারিয়েছ তার জন্য আফসোস করো না, আর যা পেয়েছ তাতে দম্ভ করো না। কারণ সব আল্লাহর \textbf{\lat{plan} (পরিকল্পনা)}। যা পেয়েছ তা আল্লাহর \textbf{\lat{gift} (উপহার)} — নিজের কৃতিত্ব নয়; তাই দম্ভ (উজব) করা আল্লাহ অপছন্দ।

\subsubsection*{৩. মূল আরবি}
\begin{center}\arab{لِّكَيْلَا تَأْسَوْا عَلَىٰ مَا فَاتَكُمْ وَلَا تَفْرَحُوا بِمَا آتَاكُمْ ۗ وَاللَّهُ لَا يُحِبُّ كُلَّ مُخْتَالٍ فَخُورٍ}\end{center}
\textbf{উচ্চারণ:} লিকাইলা- তা'সাও 'আলা- মা- ফা-তাকুম ওয়া লা- তাফরাহূ বিমা- আ-তাকুম, ওয়াল্লা-হু লা- ইউহিব্বু কুল্লা মুখতা-লিন ফাখূর।

\subsubsection*{৪. বাংলা অর্থ (\lat{pure})}
\textbf{\lat{“}যাতে তোমরা যা হারিয়েছ তার জন্য আফসোস না করো, আর যা তোমাদের দিয়েছ তাতে গর্ব (দম্ভ) না করো। আর আল্লাহ কোনো দাম্ভিক, আত্মপ্রশংসাকারীকে পছন্দ করেন না।\lat{”}}

\subsubsection*{৫. শব্দের ব্যাখ্যা (\lat{Vocabulary})}

\begin{table}[h]\centering\small
\begin{tabular}{ll}
শব্দ & অর্থ \\
\hline
\textbf{\arab{لِّكَيْلَا} \arab{تَأْسَوْا}} & যাতে তোমরা আফসোস না করো \\
\textbf{\arab{فَاتَكُمْ}} & যা তোমাদের অতিক্রম করেছে — হারিয়েছ \\
\textbf{\arab{آتَاكُمْ}} & যা তোমাদের দেওয়া হয়েছে \\
\textbf{\arab{مُخْتَالٍ} \arab{فَخُورٍ}} & দাম্ভিক, আত্মপ্রশংসাকারী \\
\hline
\end{tabular}
\end{table}

\subsubsection*{৬. গভীর তাফসীর (\lat{Deep} \lat{Tafsir})}

\begin{itemize}
\item \textbf{ইবনে কাসীর:} আয়াতটি মানুষকে দুটি জিনিস থেকে বিরত রাখে: (১) হারানো জিনিসের জন্য মাত্রাতিরিক্ত আফসোস, (২) প্রাপ্ত নেয়ামতে ফখর করা। দ্বিতীয়টি উজব ও কিবরের মূল। (ইবনে কাসীর, ৪/৩৬৯)
\item \textbf{তাবারী:} যা দুনিয়ায় মিলেছে তা আল্লাহর বণ্টন, নিজের যোগ্যতা নয় — এটি মনে না রাখলে মানুষ মুখতাল ও ফাখুর হয়ে যায়। (তাবারী, ২৩/১৫০)
\item \textbf{মাআরিফুল কুরআন:} নিজের প্রাপ্তিকে নিজের কৃতিত্ব মনে করা — এটাই উজব। আয়াত বলছে: যা পেয়েছ তাতে দম্ভ নয়; যা হারিয়েছ তাতে হতাশা নয়। এ দুটিই ঈমানের বিপরীত মনোভাব। (মাআরিফুল কুরআন, ৯/৪৬৩)
\end{itemize}
\subsubsection*{৭. সংশ্লিষ্ট হাদিস ও আয়াত}

\begin{itemize}
\item \textbf{আয়াত (নিসা ৪:৩৬):} \lat{“}নিশ্চয়ই আল্লাহ দাম্ভিক ও আত্মপ্রশংসাকারীকে পছন্দ করেন না।\lat{”}
\item \textbf{আয়াত (কাসাস ২৮:৮৩):} আখেরাত তাদের, যারা উচ্চতা চায় না।
\item \textbf{হাদিস (মুসলিম ৯১):} \lat{“}যার হৃদয়ে সরিষার দানার কিবর থাকবে, সে জান্নাতে প্রবেশ করবে না।\lat{”}
\end{itemize}
\subsubsection*{৮. গভীর শিক্ষা (\lat{Deep} \lat{Lessons})}

\begin{enumerate}
\item \textbf{প্রাপ্তি — আল্লাহর \lat{gift}।} সম্পদ, জ্ঞান, সন্তান, সফলতা — এ সব আল্লাহর \textbf{\lat{gift} (উপহার)}; নিজের যোগ্যতা মনে করা উজব।
\item \textbf{হারানো নিয়ে আফসোস নিষিদ্ধ।} যা ঘটেছে তা আল্লাহর পরিকল্পনায় ছিল — আফসোস তাকদিরের উপর \textbf{\lat{complaint} (অভিযোগ)}।
\item \textbf{দম্ভ — উজবের মূল।} সফলতা পেয়ে নিজেকে বড় ভাবা — এটাই উজব; আয়াতটি একে সরাসরি নিষিদ্ধ করেছে।
\item \textbf{ভারসাম্যপূর্ণ জীবন।} মুমিন না আফসোসে ডুবে, না গর্বে উড়ে — \textbf{\lat{balanced} (ভারসাম্যপূর্ণ)}।
\end{enumerate}
\subsubsection*{৯. জীবনে প্রয়োগ (\lat{Practical} \lat{Application})}

\begin{enumerate}
\item সফলতার পর নিজের কৃতিত্বের প্রশংসা না করে আল্লাহর \textbf{\lat{gratitude} (কৃতজ্ঞতা)} প্রকাশ করুন।
\item হারানো জিনিসের জন্য দীর্ঘ আফসোস করবেন না — দোয়া করুন ধৈর্যের জন্য।
\item নিজের সম্পদ-সন্তান-জ্ঞানকে আল্লাহর \textbf{\lat{trust} (আমানত)} মনে করুন — দম্ভ নয়।
\end{enumerate}
\medskip\hrule\medskip

\subsection*{সারসংক্ষেপ (\lat{Summary})}

\begin{table}[h]\centering\small
\begin{tabular}{lll}
আয়াত & মূল বিষয় & শিক্ষা \\
\hline
লুকমান ৩১:১৮ & অহংকারে মুখ ফেরানো ও দম্ভে হাঁটা নিষিদ্ধ & বিনয়ই মর্যাদা \\
নাহল ১৬:২৩ & আল্লাহ অহংকারীদের পছন্দ করেন না & অহংকার আল্লাহর ভালোবাসার বাধা \\
গাফির ৪০:৫৬ & বিনা দলিলে তর্কের মূল কারণ কিবর & অহংকারের লক্ষ্য কখনো অর্জিত হয় না \\
যুমার ৩৯:৭২ & অহংকারীদের আবাস জাহান্নাম & অহংকারের পরিণতি স্থায়ী \\
বনী ইসরাঈল ১৭:৩৭ & মানুষের দুর্বলতা — দম্ভ নিরর্থক & শক্তি সীমিত, দম্ভ নির্বুদ্ধিতা \\
আ'রাফ ৭:৪০ & অহংকারী জান্নাতে প্রবেশ করবে না (উট-সূঁচের মিসাল) & অহংকার জান্নাতের পথ বন্ধ করে \\
সাফফাত ৩৭:৩৫ & কালেমা শুনে অহংকার — কিবরের চূড়া & অহংকার ঈমানের প্রতিবন্ধক \\
বাকারাহ ২:২০৬ & নসীহত শুনে অহংকার — কিবরের লক্ষণ & উপদেশ মেনে নেওয়া বিনয়ের চিহ্ন \\
আ'রাফ ৭:১৪৬ & অহংকারীর অন্তর থেকে সত্যবোধ কেড়ে নেওয়া হয় & হেদায়াত হারানো সবচেয়ে বড় শাস্তি \\
কাসাস ২৮:৮৩ & আখেরাত তাদের, যারা উচ্চতা ও ফাসাদ চায় না & বিনয় আখেরাতের চাবিকাঠি \\
নিসা ৪:৩৬ & দাম্ভিক-আত্মপ্রশংসাকারী আল্লাহর অপছন্দ & ভালো কাজেও অহংকার আসতে পারে \\
হাদীদ ৫৭:২৩ & প্রাপ্তিতে দম্ভ (উজব) নিষিদ্ধ & সব নেয়ামত আল্লাহর \lat{gift} \\
\hline
\end{tabular}
\end{table}

\medskip\hrule\medskip

\subsection*{উৎসগ্রন্থ তালিকা}

\begin{table}[h]\centering\small
\begin{tabular}{ll}
গ্রন্থ & ব্যাপ্তি \\
\hline
তাফসীর ইবনে কাসীর & ১২টি আয়াতের ব্যাখ্যা \\
জামিউল বায়ান (তাবারী) & মুখ ফেরানো, তর্ক, তাকদির সংক্রান্ত ব্যাখ্যা \\
তাফসীরুল কুরতুবী & দম্ভে হাঁটা, কিবরের প্রকার \\
তাফসীর আস-সা'দী & কিবর-উজবের সংজ্ঞা ও পরিণতি \\
মাআলিমুত তানযীল (বাগাবী) & গোপন-প্রকাশ্য আমল, কালেমার অহংকার \\
মাআরিফুল কুরআন (মুফতী মুহাম্মদ শাফী উসমানী) & প্রেক্ষাপট ও শিক্ষা \\
ফী যিলালিল কুরআন (সাইয়িদ কুতুব) & আধুনিক দৃষ্টিভঙ্গি \\
\hline
\end{tabular}
\end{table}


\end{document}
